% concatenated from paper.tex
\documentclass{amsart}

%PACKAGES
%%These 3 are for editing purposes; delete before submitting
\usepackage{color}
\usepackage{soul}
\usepackage[pagewise]{lineno}
%\linenumbers

%%hyperlinks for references
\usepackage{hyperref}
\hypersetup{
    colorlinks=true, %set true if you want colored links
    linktoc=all,     %set to all if you want both sections and subsections linked
    linkcolor=blue,  %choose some color if you want links to stand out
}

%MACROS
\newcommand{\wrong}[1]{\textbf{#1}}
%%Standard math classes
\newcommand{\Q}{\mathbb{Q}}
\newcommand{\C}{\mathbb{C}}
\newcommand{\D}{\mathbb{D}}
\newcommand{\Z}{\mathbb{Z}}
\newcommand{\R}{\mathbb{R}}
\newcommand{\N}{\mathbb{N}}

%%Standard math operations
\newcommand{\dom}{\operatorname{dom}}
\newcommand{\ran}{\operatorname{ran}}
\newcommand{\diam}{\operatorname{diam}}
\newcommand{\norm}[1]{\left\| #1 \right\|}
\newcommand{\abs}[1]{\left\vert #1 \right\vert}
\newcommand{\Proj}{\operatorname{Proj}}
\newcommand{\unit}{\mathbf{1}}
\renewcommand{\star}{\ast}

%%Other math abbreviations
\newcommand{\B}{\mathcal{B}}
\renewcommand{\H}{\mathcal{H}}
\newcommand{\boldA}{\mathbf{A}}
\newcommand{\boldB}{\mathbf{B}}
\newcommand{\boldC}{\mathbf{C}}
\newcommand{\boldF}{\mathbf{F}}
\newcommand{\boldG}{\mathbf{G}}
\newcommand{\zerovec}{\mathbf{0}}
\newcommand{\zeromat}{\mathbf{0}}
\newcommand{\mvn}{\sim_{\operatorname{mvn}}}
\newcommand{\boldP}{\mathbf{P}}
\newcommand{\cstar}{$\mathrm{C}^*$}
\newcommand{\BratD}{\mathcal{D}}
\newcommand{\calC}{\mathcal{C}}
\newcommand{\calE}{\mathcal{E}}
\newcommand{\calF}{\mathcal{F}}
\newcommand{\calG}{\mathcal{G}}
\newcommand{\calH}{\mathcal{H}}
\newcommand{\calK}{\mathcal{K}}
\newcommand{\calM}{\mathcal{M}}
\newcommand{\calS}{\mathcal{S}}
\newcommand{\Id}{\operatorname{Id}}
\newcommand{\Ad}{\operatorname{Ad}}
\newcommand{\embed}{\mathcal{E}}
\newcommand{\uK}[1]{K_0^{\operatorname{sc}}(#1)}
\newcommand{\uKnoarg}{K_0^{\operatorname{sc}}}
\newcommand{\Cat}{\mbox{\bf Cat}}
\newcommand{\Tot}{\mbox{\bf Tot}}
\newcommand{\frakA}{\mathfrak{A}}
\newcommand{\frakG}{\mathfrak{G}}

%%Logic macros
\newcommand{\zerodeg}{\mathbf{0}}
\newcommand{\Deg}{\operatorname{Deg}}

%ENVIRONMENTS
%%Theorem-like environments
\theoremstyle{plain}
\newtheorem{theorem}{Theorem}[section]
\newtheorem{lemma}[theorem]{Lemma}
\newtheorem{corollary}[theorem]{Corollary}
\newtheorem{proposition}[theorem]{Proposition}
\newtheorem{conjecture}[theorem]{Conjecture}
\newtheorem{question}[theorem]{Question}
\newtheorem{fact}[theorem]{Fact}
\newtheorem{maintheorem}{Main Theorem}

%%Definition-like environments
\theoremstyle{definition}
\newtheorem{definition}[theorem]{Definition}
\newtheorem{example}[theorem]{Example}

\numberwithin{equation}{section}

%Guide to abbraviations used in labels
%%Abl- abelian group
%%AF - almost finite algebra
%%Alg- cstar algebra
%%Back- Background
%%Brat- Bratteli
%%Cat- category/categorical
%%Ce- computably enumerable
%%Cert- certificate
%%Cls-classification
%%Clscl - classical
%%Cnvx- convex ordered subgroup
%%Comp- computable
%%Dim- dimension(al)
%%Eqv- equivalence
%%Fin- finite
%%Gen- gnerated
%%Grp- group
%%Hm- homomorphism
%%Ind- Inductive
%%Indx- index
%%Iso- Isomorphism
%%Kthy- K theory
%%Lbl-labeled
%%Lim- Limit
%%Mtrx- Matrix
%%Pth- path
%%Mult- multiplicity
%%Ord- ordered
%%Prlm- preliminary
%%Pres- presentation
%%Prsnt- presentable
%%Red- reduction
%%Seq-Sequence
%%Simpl- simplicial
%%Tel- telescope
%%Uni- unique
%%Unit- Unit
%%Upr- upper


\begin{document}
\title{The computable functional calculus}
\author[C. J. Eagle]{Christopher J. Eagle${}^1$} 
\address[C. J. Eagle]
{Department of Mathematics and Statistics, University of Victoria. PO BOX 1700 STN CSC, Victoria, British Columbia, Canada. V8W 2Y2}%
\email{eaglec@uvic.ca}
\urladdr{http://web.uvic.ca/~eaglec}
\thanks{${}^1$ Supported by NSERC Discovery Grant RGPIN-2021-02459}

\author[T. H. McNicholl]{Timothy H. McNicholl}
\address[T. H. McNicholl]{Iowa State University}
\email{mcnichol@iastate.edu}
\urladdr{https://faculty.sites.iastate.edu/mcnichol/}
% \thanks{${}^3$}

\begin{abstract}
We show that the continuous functional calculus is computable.  As consequences we obtain the computable compactness of the spectrum of any computable normal element of a computably presented \cstar-algebra, the existence of effective approximate units for computably presented \cstar-algebras, and an effective version of the Spectral Theorem for compact operators on separable Hilbert spaces.
\end{abstract}
\subjclass[2020]{03D78,03D45,46L80,46L35}
\maketitle

\tableofcontents

\section{Introduction}\label{sec:Intro}
Given an element $a$ of a \cstar-algebra $\boldA$ and a polynomial $f$, there is a natural way of producing another element $f(a) \in \boldA$.  The continuous functional calculus, which is a fundamental tool in the study of \cstar-algebras, extends this to the case where $f$ is no longer required to be a polynomial, but just a function that is continuous on the spectrum of $a$, provided that $a$ is a normal element.

Following work of Fox \cite{fox2022computable}, there has recently been considerable activity in the study of computable \cstar-algebras (for example, \cite{Gelfand, AFPaper, UHFPaper, NonComputabilityPaper, FoxGoldbringHart.2024, Goldbring.2026, McNicholl.2025}).  In some of the arguments in these papers there has been a need for a computable version of the continuous functional calculus; that is, given a computable normal point $a$ of a presentation $\boldA^\#$ of $\boldA$, and a computable function $f$ on the spectrum of $a$, one frequently wants to know that $f(a)$ is also a computable point of $\boldA^\#$.  So far, these situations have been handled in an ad-hoc way, for example by using a power series representation for $f$ in the case where $f$ is analytic.  In this paper we give a systematic treatment of the continuous functional calculus in a computability setting, which provides a unified way to handle arguments of this kind.  Specifically, in Theorem \ref{theorem:CompFuncCalc} and Theorem \ref{theorem:ComputableCompactSpectrum} we prove:

\begin{theorem}
Suppose that $\boldA^\#$ is a computable presentation of a unital \cstar-algebra $\boldA$ and that the unit of $\boldA$ is a computable point.  Let $a$ be a normal computable point of $\boldA$.  Then $\sigma(a)$, the spectrum of $a$, is a computably compact subset of $\mathbb{C}$ and the continuous functional calculus is a computable map from $C(\sigma(a))$ to $\boldA^\#$.
\end{theorem}

In particular, this shows that the application of a computable function on $\sigma(a)$ to $a$ produces a computable point of $\boldA^\#$ (Corollary \ref{corollary:ApplyingFuncCalcUnital}).  We also extend our results to the non-unital setting.  Some of our results generalize work of Brattka and Dillage \cite{BrattkaDillhage.2005} and of Dillhage \cite{Dillhage.2008} from the case of a single operator on a Hilbert space to the setting of \cstar-algebras.  As a sample application, we show that if $\boldA^\#$ is a computable presentation of a \cstar-algebra $\boldA$, then there is a computable sequence of $\boldA^\#$ that is an effective approximate unit for $\boldA^\#$ (Theorem \ref{thm:ApproximateUnit}).

In Section \ref{section:CompactOperators}, we study compact operators on Hilbert spaces.  Classically, compact operators on Hilbert spaces are equivalently characterized as either those operators where the image of the closed unit ball is pre-compact or those operators that can be uniformly approximated by finite-rank operators.  We describe natural computablity-theoretic versions of both properties and prove (Theorem \ref{theorem:EquivalentCompactness}) that the two are equivalent.  We conclude by applying our computable functional calculus to the \cstar-algebra of compact operators to obtain an effective version of the Spectral Theorem (Theorem \ref{thm:Spectral}).

\section{Preliminaries}\label{}
Throughout this paper, all \cstar-algebras are assumed to be separable.  The setting we use to study the computability of \cstar-algebras, Hilbert spaces, and related structures is the one presented in \cite{EMST}, and we refer the reader there for a detailed introduction.  The necessary background can also be found in \cite{UHFPaper}.

When we consider the \cstar-algebra $\mathbb{C}$, we always view it as having its standard presentation where the only special point is $1$ (and thus the generated points are $\mathbb{Q}(i)$).  We identify $\mathbb{C}$ with this presentation throughout the paper.  To avoid repetition, \emph{throughout the paper $\boldA$ is a \cstar-algebra and $\boldA^\#$ is a fixed computable presentation of $\boldA$}.  

Our \cstar-algebra $\boldA$ may or may not be unital.  Even if it is unital, we do not have a guarantee that the unit will be a computable point of $\boldA^\#$.

\begin{definition}
Suppose that $\boldA$ is unital.  We say that $\boldA^\#$ is \emph{computably unital} if the unit $1_\boldA$ is a computable point of $\boldA^\#$.
\end{definition}

We note that every presentation of a commutative unital \cstar-algebra is computably unital \cite[Theorem 4.1]{Gelfand}, but the proof is necessarily non-uniform \cite[Theorem 5.3]{McNicholl.2025}.  More generally, every presentation of a stably finite \cstar-algebra is computably unital \cite[Proposition 4.5]{UHFPaper}.

Recall that the \emph{unitization} of a \cstar-algebra $\boldA$ (where $\boldA$ may or may not already be unital) is the unique \cstar-algebra $\boldA_1$ such that $\boldA$ is a closed two-sided ideal of $\boldA_1$ and $\boldA_1/\boldA \cong \mathbb{C}$.  Concretely, $\boldA_1 = \boldA \oplus \mathbb{C}$ as a vector space, with multiplication $(a, \lambda)\cdot(b, \mu) = (ab+\lambda b + \mu a, \lambda\mu)$ and involution $(a, \lambda)^* = (a^*, \overline{\lambda})$.  The unique norm making this *-algebra into a \cstar-algebra is given by
\[\norm{(a, \lambda)}_{\boldA_1} = \sup\{\norm{ab+\lambda b}_{\boldA} : b \in \boldA, \norm{b}_\boldA \leq 1\}.\]
The unit of $\boldA_1$ is $(0, 1)$.

The following fact was proved in \cite[Proposition 2.21]{UHFPaper}.

\begin{fact}[Computable unitization]\label{prop:CompUnit}
There is a presentation $\boldA^\#_1$ of the unitization $\boldA_1$ of $\boldA$ such that $\boldA^\#_1$ is computable and computably unital and the inclusion map is a computable map from $\boldA^\#$ to $\boldA_1^\#$.  Furthermore, an index of $\boldA^\#_1$ can be computed from an index of $\boldA^\#$.
\end{fact}

\begin{definition}
Suppose that $\boldA$ is unital. The \emph{spectrum} of an element $a \in \boldA$ is
\[\sigma_\boldA(a) = \{\lambda \in \mathbb{C} : a - \lambda1_\boldA \text{ is not invertible}\}.\]
If $\boldA$ is non-unital, then by definition $\sigma_\boldA(a) = \sigma_{\boldA_1}(a)$.
\end{definition}

When the \cstar-algebra $\boldA$ is clear from context we omit it from the notation for the spectrum.  Note, however, that if $\boldA$ is unital then it may be the case that $\sigma_{\boldA}(a) \neq \sigma_{\boldA_1}(a)$, because the unit of $\boldA$ is not the same as the unit of $\boldA_1$.  In fact, if $a \in \boldA$ then $a$ cannot be invertible in $\boldA_1$, so $0 \in \sigma_{\boldA_1}(a)$, regardless of whether or not $0 \in \sigma_\boldA(a)$.  This is the only possible difference between the spectra, so the relationship between $\sigma_{\boldA}(a)$ and $\sigma_{\boldA_1}(a)$ when $\boldA$ is unital is that $\sigma_{\boldA_1}(a) = \sigma_{\boldA}(a) \cup \{0\}$; see \cite[paragraph after Remark 2.2.1]{Murphy.1990}

In Section \ref{section:CompactOperators} we will consider operators on Hilbert spaces.  The following well-known fact was shown, in slightly different terminology, in \cite[Lemma 7]{PourElRichards.1989}; for the version in the terminology we use in this paper, see \cite[Lemma 3.1]{BrattkaYoshikawa}.

\begin{fact}\label{fact:ONB}
Every computable presentation of a Hilbert space has a computable orthonormal basis.
\end{fact}

We also note that an orthonormal basis is always c.e. closed.  This follows immediately from a more general fact about computable sequences in metric spaces.

\begin{lemma}\label{lem:ONB.ce.closed}
Let $X^\#$ be a presentation of a metric space $X$ and let $(x_n)_{n \in \mathbb{N}}$ be a computable sequence of $X^\#$ such that $\{x_n : n \in \mathbb{N}\}$ is closed and discrete.  Then $\{x_n : n \in \mathbb{N}\}$ is c.e. closed.  In particular, every computable orthonormal basis of a Hilbert space is a c.e. closed set.
\end{lemma}

\begin{proof}
Let $C = \{x_n : n \in \mathbb{N}\}$.  We describe an algorithm for enumerating the rational open balls that intersect $C$.

Enumerate the rational open balls of $X^\#$ as $(B(a_k; r_k))_{k \in \mathbb{N}}$.  At stage $s$, for each $i, j \leq s$, calculate a rational approximation to the distance between $x_i$ and $a_j$ to within an error of $2^{-s}$, say the result is $d^s_{i,j}$.  Since $d^s_{i, j}$, $2^{-s}$, and $r_j$ are all rational, we can test if $d^s_{i,j}+2^{-s} < r_j$.  If so, add $B(a_j; r_j)$ to the output.  Note that if $x_i \in B(a_j; r_j)$ for some $i$ and $j$ then for sufficiently large $s$ we will have $r_j - d(x_i, a_j) < 2^{-s}$, so we will have $d^s_{i,j}+2^{-s} < r_j$ and $B(a_j; r_j)$ will be enumerated.  On the other hand, if $x_i \not\in B(a_j; r_j)$ for all $i$ then we will never have $d^s_{i,j}+2^{-s} < r_j$, so $B(a_j; r_j)$ will not be enumerated.

The claim about orthonormal bases follows form the fact that if $(e_n)_{n \in \mathbb{N}}$ is an orthonormal basis of a Hilbert space then $\norm{e_n - e_m} = \sqrt{2}$ for all $n \neq m$.
\end{proof}

\section{Continuous functional calculus}
Recall that an element of $\boldA$ is \emph{normal} if $aa^* = a^*a$.  For normal elements of \cstar-algebras we have the continuous functional calculus:

\begin{fact}[Continuous functional calculus]
Suppose that $\boldA$ is unital.  Given a normal element $a \in \boldA$, there is a unique *-homomorphism $\Phi : C(\sigma(a)) \to \boldA$ satisfying $\Phi(1) = 1_\boldA$ and $\Phi(z) = a$, where $1$ is the constant function $1$ and $z : \sigma(a) \to \mathbb{C}$ is the identity function.  The map $\Phi$ is a *-isomorphism from $C(\sigma(a))$ onto $C^*(1_\boldA, a)$.
\end{fact}

We refer the reader to \cite[Section 2.1]{Murphy.1990} for an exposition of the continuous functional calculus and a proof of the fact above.  Given a continuous function $f : \sigma(a) \to \mathbb{C}$, we often write $f(a)$ instead of $\Phi(f)$.

Our goal in this section is to obtain a computable version of the continuous functional calculus and some of its consequences.  

\begin{definition}
Given a normal element $a$ in a unital \cstar-algebra $\boldA$, the \emph{standard presentation} of $C^*(1_\boldA, a)$ is the presentation that has $1_\boldA$ and $a$ as its special points.

Given a compact set $K \subseteq \mathbb{C}$, the \emph{standard presentation} of $C(K)$ is the presentation whose special points are the constant function $1$ and the identity function $z : K \to \mathbb{C}$.
\end{definition}

It follows from the Weierstrass approximation theorem that the standard presentation of $C(K)$ is, in fact, a presentation.  Whenever we consider computability properties of $C^*(1_\boldA, a)$ or $C(\sigma(a))$, we use these standard presentations unless we explicitly state otherwise.

\begin{theorem}\label{theorem:CompFuncCalc}
Suppose that $\boldA$ is unital, $\boldA^\#$ is computably unital, and that $a \in \boldA$ is a normal computable point of $\boldA^\#$.  Let $\Phi : C(\sigma(a)) \to C^*(1_\boldA, a)$ be the continuous functional calculus map.  Then:
\begin{enumerate}
\item{$\Phi$ is a computable *-isomorphism from the standard presentation of $C(\sigma(a))$ to the standard presentation of $C^*(1_\boldA, a)$.}
\item{The standard presentation of $C(\sigma(a))$ is computable.}
\item{$\Phi$ is a computable map from the standard presentation of $C(\sigma(a))$ to $\boldA^\#$.}
\end{enumerate}
\end{theorem}
\begin{proof}
By definition of the standard presentations of $C(\sigma(a))$ and $C^*(1_\boldA, a)$, the map $\Phi$ is computable on the generated points of $C(\sigma(a))$.  Since $\Phi$ is a *-isomorphism, this suffices to allow us to conclude that $\Phi$ is a computable *-isomorphism, establishing (1).

Since $\boldA^\#$ is computably unital, $1_\boldA$ is a computable point of $\boldA^\#$, and by hypothesis so is $a$.  Thus, since $\boldA^\#$ is computable, the standard presentation of $C^*(1_\boldA, a)$ is computable.  By (1), the standard presentation of $C(\sigma(a))$ is computably *-isomorphic to the standard presentation of $C^*(1_\boldA, a)$, which proves (2).

For statement (3), because $1_\boldA$ and $a$ are computable points of $\boldA^\#$, the inclusion map from $C^*(1_\boldA, a)$ to $\boldA$ is a computable map from the standard presentation of $C^*(1_\boldA, a)$ to $\boldA^\#$, and thus (3) follows from (1).
\end{proof}

In order to use Theorem \ref{theorem:CompFuncCalc} to apply standard functional calculus arguments in an effective setting, we first need to to establish the computable compactness of the spectrum of a normal element.  We will make use of the following definition from \cite{McNicholl.2025}.

\begin{definition}
Let $X$ be a compact Polish space and let $X^\#$ be a presentation of $X$.  A presentation $C(X)^\#$ of $C(X)$ is \emph{evaluative over $X^\#$} if the evaluation map $C(X) \times X \to \mathbb{C}$ given by $(f,x) \mapsto f(x)$ is a computable map from the induced presentation on $C(X) \times X$ to $\mathbb{C}$.
\end{definition}

\begin{theorem}\label{theorem:ComputableCompactSpectrum}
Suppose that $\boldA$ is unital, $\boldA^\#$ is computably unital, and that $a \in \boldA$ is a normal computable point of $\boldA^\#$.  Then $\sigma(a)$ is a computably compact subset of (the standard presentation of) $\mathbb{C}$.
\end{theorem}
\begin{proof}
Since the standard presentation of $C(\sigma(a))$ is computable by Theorem \ref{theorem:CompFuncCalc}, it follows from \cite[Theorem 1.1 and 1.2]{McNicholl.2025} that there is a computably compact presentation $\sigma(a)^\#$ of $\sigma(a)$ over which $C(\sigma(a))^\#$ is evaluative.  The identity function $z : \sigma(a) \to \mathbb{C}$ is a special point of the standard presentation of $C(\sigma(a))$, so since this presentation is evaluative over $\sigma(a)^\#$, the identity map is computable from $\sigma(a)^\#$ to $\mathbb{C}$.  As $\sigma(a)^\#$ is computably compact and computable compactness is preserved by computable maps (see, e.g., \cite[Lemma 3.31]{DowneyMelnikov}), we have that $\sigma(a)$ is a computably compact subset of $\mathbb{C}$.
\end{proof}

We can extend this to the case where $\boldA$ is not necessarily unital and where $\boldA$ is unital but $\boldA^\#$ is not computably unital.

\begin{theorem}\label{theorem:ComputableCompactSpectrumNonunital}
Let $\boldA^\#$ be a computable presentation of a (not necessarily unital) \cstar-algebra $\boldA$.  If $a$ is a normal computable point of $\boldA^\#$ then $\sigma(a)$ is a computably compact subset of $\mathbb{C}$.
\end{theorem}
\begin{proof}
Let $\boldA_1$ be the unitization of $\boldA$.  By Fact \ref{prop:CompUnit} there is a computably unital computable presentation $\boldA_1^\#$ of $\boldA_1$.  The inclusion map from $\boldA$ to $\boldA_1$ is computable, so if $a$ is a computable point of $\boldA^\#$ then it is also a computable point of $\boldA_1^\#$.  Theorem \ref{theorem:ComputableCompactSpectrum} implies that $\sigma_{\boldA_1}(a)$ is a computably compact subset of $\mathbb{C}$.  

If $\boldA$ is non-unital, then $\sigma_{\boldA}(a) = \sigma_{\boldA_1}(a)$ by definition, and so $\sigma_{\boldA}(a)$ is computably compact.

If $\boldA$ is unital, then we have $\sigma_{\boldA_1}(a) = \sigma_{\boldA}(a) \cup \{0\}$ (see \cite[paragraph after Remark 2.2.1]{Murphy.1990}).  

If $a$ is non-invertible in $\boldA$ then $0 \in \sigma_\boldA(a)$, so $\sigma_{\boldA_1}(a) = \sigma_{\boldA}(a)$ and $\sigma_{\boldA}(a)$ is thus computably compact.

If $a$ is invertible in $\boldA$ then $0 \not\in \sigma_{\boldA}(a)$, and in fact the distance from $0$ to $\sigma_{\boldA}(a)$ is precisely $1/r(a^{-1})$, where $r$ is the spectral radius.  Since $a$ is normal we have $r(a^{-1}) = \norm{a^{-1}}$, so the distance from $0$ to $\sigma_\boldA(a)$ is $1/\norm{a^{-1}}$; it follows from the fact that $a$ is a computable point that we can compute $\norm{a^{-1}}$, so the distance from $0$ to $\sigma_{\boldA}(a)$ is a computable real number.  Given $n \in \mathbb{N}$, we can find $N \geq n$ such that $2^{-N} < 1/\norm{a^{-1}}$.  Using the computable compactness of $\sigma_{\boldA_1}(a)$ we can effectively find a finite tuple of rational open balls of radius $2^{-N}$ that cover $\sigma_{\boldA_1}(a)$.  Removing any of these balls that contains $0$ produces a cover of $\sigma_{\boldA}(a)$.
\end{proof}

We note that the proof of Theorem \ref{theorem:ComputableCompactSpectrum} is uniform.  The proof of Theorem \ref{theorem:ComputableCompactSpectrumNonunital} has two sources of non-uniformity (namely, whether or not the algebra is unital, and if so, whether or not the element is invertible).  We do not know if the non-uniformity is necessary.  However, the proof is uniform if we know that the algebra is non-unital, and this (together with the uniformity of Theorem \ref{theorem:ComputableCompactSpectrum} for the computably unital case) is sufficient uniformity for all applications of which we are aware.

Theorem \ref{theorem:CompFuncCalc} allows us to apply standard functional calculus arguments in a computability setting.  In what follows, when we discuss computable functions on $\sigma(a)$ we are thinking of $\sigma(a)$ as a computably compact subset of $\mathbb{C}$.

\begin{corollary}\label{corollary:ApplyingFuncCalcUnital}
Let $\boldA^\#$ be a computably unital presentation of a \cstar-algebra $\boldA$ and let $a$ be a normal computable point of $\boldA^\#$.  Suppose that $f : \sigma(a) \to \mathbb{C}$ is a computable function.  Then $f(a)$ is a computable point of $\boldA^\#$.
\end{corollary}
\begin{proof}
By the Effective Weierstrass Theorem \cite[Theorem I.0.6]{PourElRichards.1989}, since $f$ is a computable function with computably compact domain, $f$ is effectively uniformly approximable by polynomials, which means that $f$ is a computable point of the standard presentation of $C(\sigma(a))$.  By Theorem \ref{theorem:CompFuncCalc}, the continuous functional calculus map $\Phi : C(\sigma(a)) \to \boldA$ is computable, so $f(a) = \Phi(f)$ is a computable point of $\boldA^\#$.
\end{proof}

\begin{corollary}\label{corollary:ApplyingFuncCalc}
Let $a$ be a normal computable point of $\boldA^\#$.  Suppose that $f : \sigma(a) \cup \{0\} \to \mathbb{C}$ is a computable function and that $f(0) = 0$.  Then $f(a)$ is a computable point of $\boldA^\#$.
\end{corollary}
\begin{proof}
By Corollary \ref{corollary:ApplyingFuncCalcUnital}, $f(a)$ is a computable point of $\boldA_1^\#$.  The condition $f(0) = 0$ is equivalent having $f(a) \in \boldA$.  Since the inclusion map from $\boldA$ to $\boldA_1$ is an isometry, it follows that $f(a)$ is also a computable point of $\boldA^\#$.
\end{proof}

For later reference, we record some specific instances of Corollary \ref{corollary:ApplyingFuncCalc}

\begin{corollary}\label{corollary:PositiveNegativeParts}
If $a \in \boldA$ is a normal computable point of $\boldA^\#$ then the real and imaginary parts $a_r$ and $a_i$ of $a$ are computable points of $\boldA^\#$.  If $a$ is a self-adjoint computable point of $\boldA^\#$ then the positive and negative parts $a^+$, and $a^-$ are computable points of $\boldA^\#$.
\end{corollary}
\begin{proof}
The function $\operatorname{Re}(z)$ is computable on any computably compact subset of $\mathbb{C}$, and $\operatorname{Re}(0) = 0$, so we may apply Corollary \ref{corollary:ApplyingFuncCalc} to conclude that $a_r = \operatorname{Re}(a)$ is a computable point of $\boldA^\#$.  By similar reasoning, $a_i = \operatorname{Im}(a)$ is also a computable point of $\boldA^\#$.

If $a$ is positive then $\sigma(a) \subseteq [0, \infty)$.  The function $\sqrt{z}$ is computable on any computably compact subset of $[0, \infty)$, and $\sqrt{0} = 0$, so again we may apply Corollary \ref{corollary:ApplyingFuncCalc} to conclude that $a^{1/2}$ is a computable point of $\boldA^\#$.  

The remaining claims follow because $a^2$ is positive for any self-adjoint $a$ and we have $\abs{a} = (a^2)^{1/2}$, $a^+ = \frac{1}{2}(\abs{a}+a)$, and $a^- = \frac{1}{2}(\abs{a}-a)$.
\end{proof}

We note that the Corollaries above are uniform for non-unital algebras and also are uniform for computably unital presentations.

\section{Effective approximate units}
In many treatments of the theory of \cstar-algebras one of the early uses of the functional calculus is to show that every separable \cstar-algebra has a \emph{sequential approximate unit}.  As an example of how the results of the previous section can be applied, we show that computable \cstar-algebras have computable sequential approximate units.

\begin{definition}
A \emph{sequential approximate unit} for $\boldA$ is a sequence $(u_n)_{n \in \mathbb{N}}$ of positive elements of $\boldA$ such that for every $a \in \boldA$,
\[\lim_{n\to\infty}u_na = \lim_{n\to\infty}au_n = a.\]
\end{definition}

If $\boldA$ is unital then the constant sequence $u_n = 1_\boldA$ is an approximate unit, so the interesting case is when $\boldA$ is non-unital.  Therefore, \emph{for this section, $\boldA$ is assumed to be non-unital}. 

\begin{definition}
A sequential approximate unit is an \emph{effective approximate unit} of $\boldA^\#$ if it is a computable sequence and from a (code of a) generated point $a$ of $\boldA^\#$ it is possible to compute a modulus of convergence for $(u_na)_{n \in \mathbb{N}}$ and a modulus of convergence for $(au_n)_{n\in\mathbb{N}}$.
\end{definition}

Recall that an element $a \in \boldA$ is \emph{positive}, written $a \geq 0$, if $a$ is normal and $\sigma(a) \subseteq [0, \infty)$.  It is well-known that the positive elements of a \cstar-algebra can be equivalently described as those $a$ for which there exists a $b$ so that $a = b^*b$.  We let $\boldA^+$ denote the set of positive elements of $\boldA$.  The notion of positivity induces a translation-invariant partial order on the self-adjoint elements of $\boldA$ defined by $a \leq b$ if and only if $b - a \geq 0$.  See \cite[Section 2.2]{Murphy.1990} for a full discussion of positive elements of \cstar-algebras.

For the remainder of the section, let $\Lambda = \{a \in \boldA^+ : \norm{a} < 1\}$.  Classically, $\Lambda$ forms a net that is an approximate unit for $\boldA$, even when $\boldA$ may be non-separable (see \cite[Theorem 3.1.1]{Murphy.1990}).

\begin{lemma}\label{lemma:ComputableApproxUnitDirected}
If $a, b \in \Lambda$ are computable points of $\boldA^\#$ then $\Lambda$ contains a computable point $c$ of $\boldA^\#$ such that $a \leq c$ and $b \leq c$.  Moreover, an $\boldA^\#$-index of $c$ can be computed from $\boldA^\#$-indices of $a$ and $b$.
\end{lemma}
\begin{proof}
Working in $\boldA_1$, let
\[c = (a(1_{\boldA_1}-a)^{-1}+b(1_{\boldA_1}-b)^{-1})(1_{\boldA_1}+a(1_{\boldA_1}-a)^{-1}+b(1_{\boldA_1}-b)^{-1})^{-1}.\]
It is shown on \cite[Page 78]{Murphy.1990} that $c \in \Lambda$ and that $a \leq c$ and $b \leq c$.  Since $\norm{a} < 1$ we have $(1_{\boldA_1}-a)^{-1} = \sum_{n=0}^{\infty}a^n$, and likewise for $b$.  Since $a$ and $b$ are computable points of $\boldA^\#$, this implies that $c$ is as well.
\end{proof}

\begin{lemma}\label{lemma:ComputableApproxUnitOneStep}
Suppose that $a \in \Lambda$, $a$ is a computable point of $\boldA^\#$, and $n \in \mathbb{N}$.  Then $\Lambda$ contains a computable point $e$ of $\boldA^\#$ such that whenever $c \in \Lambda$ and $c \geq e$ we have $\max\{\norm{a-ca}, \norm{a-ac}\} < 2^{-n}$.  Moreover, an $\boldA^\#$-index of $e$ can be computed from $n$ and an $\boldA^\#$-index of $a$.
\end{lemma}
\begin{proof}
Define $h : [0, \infty) \to \mathbb{C}$ by
\[h(t) = \begin{cases}1 & t \geq 2^{-n} \\ 2^{n+1}t-1 & 2^{-(n+1)} \leq t < 2^n \\ 0 & t < 2^{-(n+1)}\end{cases}\]
Note that $h$ is continuous and $h(0) = 0$.  Let $g = (1-2^{-(n+1)})h\vert_{\sigma(a) \cup \{0\}}$.  Set $e = g(a)$.  By Corollary \ref{corollary:ApplyingFuncCalc}, $e$ is a computable point of $\boldA^\#$.  Since $h \geq 0$, it follows that $e$ is positive.  We have $\norm{e} = \norm{g} \leq 1-2^{-(n+1)} < 1$, so $e \in \Lambda$.

Suppose that $c \in \Lambda$ and $c \geq e$.  Then $\norm{a-ca} = \norm{a(1_{\boldA_1}-c} < \norm{1_{\boldA_1} - c}$.  Since $c \geq e$, $1_{\boldA_1} -c \leq 1_{\boldA_1}-e$.  However, as $\norm{c} < 1$, $1_{\boldA_1} - c \geq 0$.  Hence $\norm{1_{\boldA} - c} \leq \norm{1-g} < 2^{-n}$.  The proof that $\norm{a-ca}<2^{-n}$ is similar.
\end{proof}

\begin{lemma}
There is a c.e. set $\Lambda_0$ of generated points of $\boldA^\#$ that is dense in $\Lambda$.
\end{lemma}
\begin{proof}
Let $\Lambda_0 = \{c^*c : a \text{ is a generated point of $\boldA^\#$ and $\norm{c} < 1$}\}$.  Then $\Lambda_0 \subseteq \Lambda$ and $\Lambda_0$ is c.e..  Suppose that $a \in \Lambda$ and $\epsilon > 0$.  Since $a$ is positive, there is $b \in \boldA$ such that $a = b^*b$, and since $\norm{a} < 1$ we also have $\norm{b} < 1$.  There is a generated point $c$ of $\boldA^\#$ such that $\norm{c-b} < \epsilon(\norm{b}+1)^{-1}$ and so that $\norm{c} < 1$.  It follows that $\norm{a-c^*c} < \epsilon$.
\end{proof}

\begin{theorem}\label{thm:ApproximateUnit}
There is an effective approximate unit of $\boldA^\#$.
\end{theorem}
\begin{proof}
Fix a c.e. set $\Lambda_0$ as in the previous lemma.  We recursively define a sequence $(u_n)_{n\in\mathbb{N}}$.  Suppose that $u_j$ has been defined for all $j < n$.  By Lemma \ref{lemma:ComputableApproxUnitOneStep}, for each $k \in \{0, \ldots, n\}$, we can compute an $\boldA^\#$-index of an $e_{k,n} \in \Lambda_0$ so that $\max\{\norm{c-ce_{k,n}}, \norm{c-e_{k,n}c}\} < 2^{-n}$ whenever $c \in \Lambda$ is such that $c \geq e_{k,n}$.  By Lemma \ref{lemma:ComputableApproxUnitDirected} we can then compute an $\boldA^\#$-index of a $u_n \in \Lambda$ such that $u_n \geq e_{k, n}$ for all $k \leq n$ and so that $u_n \geq u_j$ for all $j < n$.  We will show that this sequence $(u_n)_{n\in\mathbb{N}}$ is an effective approximate unit.

Fix a generated point $a$ of $\boldA^\#$.  Let $b = (\norm{a}+1)^{-1}a$.  Let $b_r$ and $b_i$ be the real and imaginary parts of $b$, respectively, and let their respective positive and negative parts be $b_r^+, b_r^-, b_i^+$, and $b_i^-$.  By Corollary \ref{corollary:PositiveNegativeParts}, $b_r^+, b_r^-, b_i^+$, and $b_i^-$ are computable points of $\boldA^\#$.  Moreover, it is possible to compute an $\boldA^\#$-index for each of them from a $\boldA^\#$-index for $a$.

Fix $k \in \mathbb{N}$.  Set $k_0 = -\lceil \log_2(\norm{a} + 1) \rceil + k + 4$.  Compute $k(r, +) \in \mathbb{N}$ such that $\norm{b_r^+-u_{k(r,+)}} < 2^{-k_0}$, and likewise compute such $k(r, -)$, $k(i, +)$, and $k(i, -)$.  Let $N_0 = \max\{k(r,+), k(r,-), k(i,+), k(i, -)\}$.

Fix $n \geq N_0$.  Since $\norm{u_n} < 1$ and $N_0 \geq k(r,+)$,
\begin{align*}
\abs{\norm{b_r^+ - u_n b_r^+} - \norm{u_{k(r,+)} - u_n u_{k(r,+)}}}
&= \abs{\norm{b_r^+} - \norm{u_{k(r,+)}}}\norm{1_{\boldA_1}-u_n} \\ &\leq \norm{b_r^+-u_{k(r,+)}}(\norm{1_{\boldA_1}}+\norm{u_n}) \\
&< 2\norm{b_r^+ - u_{k(r,+)} } \\
&< 2\cdot2^{-k_0}
\end{align*}
thus
\[\norm{b_r^+ - u_nb_r^+} < \norm{u_{k(r,+)} - u_nu_{k(r,+)}} + 2\cdot2^{-k_0}.\]
Also, by construction, $\norm{u_{k(r,+)}-u_nu_{k(r,+)}} < 2^{-k_0}$, so we obtain
\[\norm{b_r^+ - u_nb_r^+} < 3\cdot2^{-k_0}.\]
Similar calculations produce the same estimate using $b_r^-$, $b_i^+$, and $b_i^-$. The decomposition $b = b_r^+ - b_r^- + ib_i^+ - ib_i^-$
therefore implies
\[
\norm{b-u_nb} < 4\cdot3\cdot2^{-k_0} < 2^{-k_0+4}.
\]
Thus, by definition of $k_0$,
\[\norm{b-u_nb} < (\norm{a}+1)2^{-k}.\]
Finally, since $b = (\norm{a}+1)^{-1}a$, this implies the desired
\[\norm{a-u_na} < 2^{-k}.\]
It follows similarly that $\norm{a - a u_n} < 2^{-k}$.
\end{proof}

\section{Spectral theory for compact operators}\label{section:CompactOperators}
We now turn our attention to compact operators on separable Hilbert spaces.  Throughout this section, let $\mathcal{H}$ be a separable Hilbert space and $\mathcal{H}^\#$ be a presentation of $\mathcal{H}$.  Let $(\rho_j)_{j \in \mathbb{N}}$ be an enumeration of the rational vectors of $\mathcal{H}^\#$.  For each non-zero $v \in \mathcal{H}$, let $P_v$ denote the orthogonal projection onto $v$.  We refer to $P_{\rho_j}$ as the \emph{$j$th rational rank $1$ projection}.  We denote by $\mathcal{K}(\mathcal{H})$ the (non-unital) \cstar-algebra of compact operators on $\mathcal{H}$.

\subsection{Presentations of the compact operators}
Our first goal is to use $\mathcal{H}^\#$ to produce a presentation of $\mathcal{K}(\mathcal{H})$.

\begin{lemma}\label{lem:ProjectionNorm}
For all $v, w \in \mathcal{H}$,
\[\norm{P_v - P_w} \leq 2\norm{\norm{v}^{-1}v - \norm{w}^{-1}w}.\]
\end{lemma}
\begin{proof}
First, suppose that $v, w, u \in \mathcal{H}$ are unit vectors.  Then we have:
\begin{align*}
\norm{P_v(u) - P_w(u)} &= \norm{\langle u,v \rangle v - \langle u, w \rangle w} \\
&= \norm{\langle u, v \rangle v - \langle u, v \rangle w + \langle u, v \rangle w - \langle u, w \rangle w} \\
&\leq \abs{\langle u, v \rangle}\norm{v-w} + \abs{\langle u,v \rangle - \langle u, w \rangle}\norm{w} \\
&= \abs{\langle u, v\rangle}\norm{v-w} + \abs{\langle u, v-w \rangle} \\
&\leq \norm{v-w} + \norm{v-w} &\text{(Cauchy-Schwartz)}\\
&= 2\norm{v-w}
\end{align*}
Thus when $\norm{v} = \norm{w} = 1$ we have $\norm{P_v - P_w} \leq 2\norm{v-w}$.  Since $P_{\lambda v} = P_v$ for all $\lambda \in \mathbb{C}$, and likewise for $w$, rescaling gives the desired conclusion.
\end{proof}

\begin{lemma}\label{lem:rationaldense}
The span of the set of rational rank $1$ projections is dense in $\mathcal{K}(\mathcal{H})$.    
\end{lemma}
\begin{proof}
By \cite[Theorems 2.4.5 and 2.4.6]{Murphy.1990}, the span of the set of all rank $1$ projections is dense in $\mathcal{K}(\mathcal{H})$.  It therefore suffices to show that every rank $1$ projection can be approximated by rational rank $1$ projections.  The rational vectors of $\mathcal{H}^{\#}$ are dense in $\mathcal{H}$, so Lemma \ref{lem:ProjectionNorm} shows that the rational rank $1$ projections are dense in the rank $1$ projections, and hence the span of the rational rank $1$ projections is dense in $\mathcal{K}(\mathcal{H})$.
\end{proof}

\begin{definition}
By $\mathcal{K}(\mathcal{H}^\#)$ we denote the presentation of $\mathcal{K}(\mathcal{H})$ whose $j$th special point is $P_{\rho_j}$.
\end{definition}

By Lemma \ref{lem:rationaldense}, $\mathcal{K}(\mathcal{H}^\#)$ is indeed a presentation of $\mathcal{K}(\mathcal{H})$.  The norm estimate given by Lemma \ref{lem:ProjectionNorm} immediately implies:

\begin{corollary}\label{cor:ComputePtProj}
If $v$ is a computable point of $\mathcal{H}^\#$ then $P_v$ is a computable point of $\mathcal{K}(\mathcal{H}^\#)$.  Moreover, from an index for $v$ we can compute an index for $P_v$.
\end{corollary}

We next aim to show that $\mathcal{K}(\mathcal{H}^\#)$ is a computable presentation whenever $\mathcal{H}^\#$ is computable.

\begin{lemma}\label{lem:IdentityMapLift}
Suppose that $\mathcal{H}^\dagger$ is a presentation of $\mathcal{H}$ such that the identity map is a computable map from $\mathcal{H}^\#$ to $\mathcal{H}^\dagger$.  Then the identity map is a computable map from $\mathcal{K}(\mathcal{H}^\#)$ to $\mathcal{K}(\mathcal{H}^\dagger)$.
\end{lemma}
\begin{proof} 
It suffices to show that from a generated point $\rho$ of $\mathcal{K}(\calH^\#)$ and $k \in \N$ it is possible to compute 
a generated point $\rho'$ of $\mathcal{K}(\calH^\dagger)$ so that $\norm{\rho - \rho'} < 2^{-k}$. 

Let $\xi_j$ denote the $j$-th rational vector of $\calH^\dagger$. 
Suppose $\rho = p(P_{\rho_0}, \ldots, P_{\rho_n})$ where $p$ is a rational $*$-polynomial with no constant term.  
From $p$, it is possible to compute a modulus of continuity $g$ for $p$ on the unit ball of $\calH^n$; that is, a function $g : \mathbb{N} \to \mathbb{N}$ such that whenever $v_0, \ldots, v_n, u_0, \ldots, u_n$ all have norm at most $1$ and $\max_j \norm{v_j - u_j} \leq 2^{-g(m)}$ we have $\norm{p(v_0, \ldots, v_n) - p(u_0, \ldots, u_n)} < 2^{-m}$. 

Since the identity map from $\mathcal{H}^\#$ to $\mathcal{H}^\dagger$ is computable, each $\xi_j$ is a computable point of $\mathcal{H}^\#$. Thus, by Lemma \ref{lem:ProjectionNorm}, we can compute $\xi_{j_0}, \ldots, \xi_{j_n}$ so that 
$\norm{P_{\rho_s} - P_{\xi_{j_s}}} < 2^{-g(k)}$ for all $s \leq n$.  
Thus,
\[\norm{\rho - p(P_{\xi_{j_0}}, \ldots, P_{\xi_{j_n}})} < 2^{-k}.\]
\end{proof}

\begin{proposition}\label{prop:PresentationOfCompacts}
If $\mathcal{H}^\#$ is computable then so is $\mathcal{K}(\mathcal{H}^\#)$.
\end{proposition}
\begin{proof}
The proofs in the infinite-dimensional case and the finite-dimensional case are essentially identical, so we use the notation of the infinite-dimensional case.

Using Fact \ref{fact:ONB}, let $(e_n)_{n \in \mathbb{N}}$ be a computable orthonormal basis for $\mathcal{H}^\#$.  

Let $\mathcal{H}^\dagger$ be the presentation of $\mathcal{H}$ whose $n$th special point is $e_n$ and let $\xi_j$ denote the $j$th rational vector of $\mathcal{H}^\dagger$.  Since $(e_n)_{n\in\mathbb{N}}$ is a computable sequence of $\mathcal{H}^\#$, the identity map is a computable map from $\mathcal{H}^\#$ to $\mathcal{H}^\dagger$.  By Lemma \ref{lem:IdentityMapLift} the identity map is a computable map from $\mathcal{K}(\mathcal{H}^\#)$ to $\mathcal{K}(\mathcal{H}^\dagger)$ and thus is a computable *-isomorphism. Therefore, it suffices to show that $\mathcal{K}(\mathcal{H}^\dagger)$ is a computable presentation.

Let $\rho$ be a generated point of $\mathcal{K}(\mathcal{H}^\dagger)$.  Then we can write $\rho = p(P_{\xi_0}, \ldots, P_{\xi_m})$ for some rational $*$-polynomial $p$ and some $m$.  From the computability of the orthonormal basis $(e_n)_{n \in \mathbb{N}}$ we can compute the coefficients of each $\xi_j$ with respect to that orthonormal basis, and from this and $p$ it is possible to compute another *-polynomial $q$ and an $N$ such that $\rho = q(P_{e_0}, \ldots, P_{e_N})$.  Since the closed unit ball of the span of $\{e_0, \ldots, e_N\}$ is a computably compact set of $\mathcal{H}^\dagger$, this expression allows us to calculate $\norm{\rho}$.
\end{proof}

All separable Hilbert spaces are computably categorical (see \cite[Corollary 3.7]{BrattkaYoshikawa}), so between any two computable presentations of separable Hilbert spaces of the same dimension there is a computable isometric isomorphism.  We now show that computable isometric isomorphisms between Hilbert spaces lift to computable *-isomorphisms on the algebras of compact operators.

\begin{lemma}\label{lem:compConjugation}
Let $\mathcal{H}_0^\#$ and $\mathcal{H}_1^\#$ be computable presentations of Hilbert spaces $\mathcal{H}_0$ and $\mathcal{H}_1$ and let $T : \mathcal{H}_0 ^\#\to \mathcal{H}_1^\#$ be a computable isometric isomorphism.  Then the map $\Phi_T : \mathcal{K}(\mathcal{H}_0) \to \mathcal{K}(\mathcal{H}_1)$ defined by $\Phi_T(S) = TST^{-1}$ is a computable *-isomorphism from $\mathcal{K}(\mathcal{H}_0^\#)$ to $\mathcal{K}(\mathcal{H}_1^\#)$.
\end{lemma}
\begin{proof}
Since $\Phi_T$ is a $*$-homomorphism, to show that $\Phi_T$ is computable it suffices to show that $\Phi_T$ is computable on the special points of $\mathcal{K}(\mathcal{H}_0^\#)$.  So we consider a special point, which has the form $P_\xi$ for some special point $\xi$ of $\mathcal{H}_0^\#$.  Then $\Phi_T(P_\xi) = P_{T(\xi)}$.  Since $T$ is computable, we can effectively approximate $T(\xi)$ arbitrarily well by rational vectors of $\mathcal{H}_1^\#$. Then it follows from Lemma \ref{lem:ProjectionNorm} that $P_{T(\xi)}$ is a computable point of $K(\mathcal{H}_1^\#)$.
\end{proof}

For each $n \in \mathbb{N}$, let $\ell^2_n$ be the $n$-dimensional Euclidean space, and let $\ell^2$ be the infinite-dimensional separable Hilbert space of square-summable sequences.  Each of these spaces has a standard computable presentation whose special points are given by the standard orthonormal basis for that space.  Using these presentations, the following is immediate from Lemma \ref{lem:compConjugation} and the computable categoricity of Hilbert spaces.

\begin{proposition}
Suppose that $\mathcal{H}^\#$ is computable.
\begin{enumerate}
\item{If $\dim(\mathcal{H}) = n$, then $\mathcal{K}(\mathcal{H}^\#)$ is computably isomorphic to $\mathcal{K}(\ell_n^2)$.}
\item{If $\mathcal{H}$ is infinite-dimensional, then $\mathcal{K}(\mathcal{H}^\#)$ is computably isomorphic to $\mathcal{K}(\ell^2)$.}
\end{enumerate}
\end{proposition}

We have thus shown that computable presentations of $\mathcal{K}(\mathcal{H})$ arising from computable presentations of $\mathcal{H}$ are all computably isomorphic.  By \cite[Theorem 4.18]{UHFPaper}, all unital UHF algebras are computably categorical.  In the infinite-dimensional case, the compact operators $\mathcal{K}(\mathcal{H})$ form a \emph{non-unital} UHF algebra.  The methods from \cite{UHFPaper} do not apply to non-unital UHF algebras. Nevertheless, \cite[Theorem 4.18]{UHFPaper} and Lemma \ref{lem:compConjugation} together suggest the possibility of a positive answer to:

\begin{question}
Is $\mathcal{K}(\mathcal{H})$ computably categorical?    
\end{question}

\subsection{Computably compact operators}
Classically, compact operators on Hilbert spaces can be equivalently characterized as limits of finite-rank operators or as operators where the closure of the image of the closed unit ball is compact.  Our definition of $\mathcal{K}(\mathcal{H}^\#)$ is based on the former description.  The latter also has a natural computability-theoretic version.

\begin{definition}
Let $T$ be an operator on $\mathcal{H}$.  We say that $T$ is a \emph{computably image-compact operator of $\mathcal{H}^\#$} if $T$ is a computable operator of $\mathcal{H}^\#$ and $\overline{T[\overline{B}(0_\mathcal{H}; 1)]}$ is a computably compact set of $\mathcal{H}^\#$.
\end{definition}

A note on terminology is in order.  In \cite{BrattkaDillhage.2007}, the term ``computably compact operator" is used in the context of Banach spaces with the approximation property to mean a computable operator that can be computably approximated by finite-rank operators; in our setting, this is notion corresponds to computable points of $\mathcal{K}(\mathcal{H}^\#)$.  The notion of computable image-compactness, defined above, makes sense even in the context of Banach spaces without the approximation property.  We remain in the context of operators on Hilbert spaces, where our next goal is to show that computably image-compact operators of $\mathcal{H}^\#$ are exactly the computable points of $K(\mathcal{H}^\#)$, thus giving a computability analog of the classical equivalence between the definitions of compact operators on Hilbert spaces.

\begin{lemma}\label{lemma:Adjoint}
Suppose that $T$ is a computable finite-rank operator of $\mathcal{H}$.  Then the adjoint $T^*$ is also computable.  Moreover, a code for $T^*$ can be computed from a code for $T$ and the rank of $T$.
\end{lemma}
\begin{proof}
Let $K$ be the rank of $T$.  Using Fact \ref{fact:ONB}, let $(e_n)_{n \in \mathbb{N}}$ be a computable orthonormal basis of $\mathcal{H}^\#$.  

We search for indices $n_1, \ldots, n_K$ such that $\{T(e_{n_1}), \ldots, T(e_{n_K})\}$ is linearly independent.  To do this, given any indices $n_1, \ldots, n_K$, recall that $\{T(e_{n_1}), \ldots, T(e_{n_K})\}$ is linearly independent if and only if the matrix $B_{n_1, \ldots, n_K}$ whose $i,j$ entry is $\langle T(e_{n_i}), T(e_{n_j}) \rangle$ has strictly positive determinant.  Fixing a computable enumeration of $\mathbb{N}^K \times \mathbb{Q}_{>0}$, for each $(n_1, \ldots, n_K, \epsilon)$ we compute the determinant of $G_{n_1, \ldots, n_K}$ to within $\epsilon$; if the result guarantees that the determinant is strictly positive, the search is complete.  Since the rank of $T$ is exactly $K$, the search will eventually terminate.

Now we have a basis $\{T(e_{n_1}), \ldots, T(e_{n_K})\}$ for the image of $T$.  Let $G = B_{n_1, \ldots, n_K}^{-1}$, and denote the entry in position $(i, j)$ of $G$ by $G_{i,j}$.  Then for any $x \in \mathcal{H}$,
\[T^*(x) = \sum_{i=1}^K\sum_{j=1}^KG_{i,j}\langle x, T(e_{n_j}) \rangle e_{n_i}.\]
\end{proof}

\begin{lemma}\label{lem:FiniteRankApproximation1}
Suppose that $T$ is a computably image-compact operator of $\mathcal{H}^\#$, and fix $k \in \mathbb{N}$ and a computable orthonormal basis $(e_n)_{n\in\mathbb{N}}$ of $\mathcal{H}^\#$.  Then we can compute $K$, $n_1, \ldots, n_K$ such that $\norm{T - \sum_{i=1}^KP_{e_n} \circ T} < 2^{-k}$.
\end{lemma}
\begin{proof}
By Lemma \ref{lem:ONB.ce.closed}, $\{e_n : n \in \mathbb{N}\}$ is c.e. closed.  Let $G = \{n \in \mathbb{N} : \operatorname{range}(P_{v_n} \circ T) \neq \{0_{\mathcal{H}}\}\}$.  Note that $G$ is a c.e. set.

Suppose $F \subseteq G$ and $F$ is finite.  We claim that $\norm{T - \sum_{n \in F}P_{e_n} \circ T}$ is a computable real number, uniformly in $F$.  To see this, let $C = \overline{T[\overline{B}(0_{\mathcal{H}}; 1)]}$.  Then
\[\norm{T - \sum_{n \in F}P_{e_n} \circ T} = \sup_{v \in C}\norm{v - \sum_{n \in F}P_{e_n}(v)}.\]
Since $e_n$ is a computable vector of $\mathcal{H}^\#$, uniformly in $n$, $\operatorname{Id}_{\mathcal{H}} - \sum_{n \in F}P_{e_n}$ is a computable operator of $\mathcal{H}^\#$, uniformly in $F$.  By hypothesis, $C$ is a computably compact set of $\mathcal{H}^\#$, so it follows that $\norm{T - \sum_{n \in F}P_{e_n} \circ T}$ is a computable real number, uniformly in $F$.

We now search for a finite set $F \subseteq G$ such that $\norm{T - \sum_{n \in F}P_{e_n} \circ T} < 2^{-k}$.  The previous paragraph shows that this search is effective and, by \cite[Corollary 4.5]{Conway.1990}, it terminates.
\end{proof}

\begin{lemma}
Suppose that $T$ is a computably image-compact operator of $\mathcal{H}^\#$.  Then the adjoint $T^*$ is a computably image-compact operator of $\mathcal{H}^\#$ as well, and a code for $T^*$ can be computed from a code for $T$.
\end{lemma}
\begin{proof}
Using Fact \ref{fact:ONB}, let $(e_n)_{n\in\mathbb{N}}$ be a computable orthonormal basis of $\mathcal{H}^\#$, and fix $k \in \mathbb{N}$.  Using Lemma \ref{lem:FiniteRankApproximation1}, find $K$, $n_1, \ldots, n_K$ be such that $\norm{T - \sum_{i=1}^{K}P_{e_i}\circ T} < 2^{-k}$.  Let $S = \sum_{i=1}^KP_{e_i} \circ T$.  Then $S$ is a computable operator of rank $K$, so by Lemma \ref{lemma:Adjoint} $S^*$ is computable, and we can compute a code for $S^*$ from $K$ (which we have already computed) and a code for $S$, which in turn we can compute from a code for $T$.  We have $\norm{T^*-S^*} = \norm{T - S} < 2^{-k}$.
\end{proof}

\begin{lemma}\label{lem:FiniteRankApproximation2}
Suppose that $T$ is a computably image-compact operator of $\mathcal{H}^\#$, and fix $k \in \mathbb{N}$.  Then there are computable scalars $\lambda_0, \ldots, \lambda_n$ and computable vectors $v_0, \ldots, v_n$ of $\mathcal{H}^\#$ so that $\norm{T - \sum_{j \leq n}\lambda_jP_{v_j}} < 2^{-k}$.
\end{lemma}
\begin{proof}
Let $(e_n)_{n \in \mathbb{N}}$ be a computable orthonormal basis for $\mathcal{H}^\#$, and use Lemma \ref{lem:FiniteRankApproximation1} to find $K$, $n_1, \ldots, n_K$ such that $\norm{T - \sum_{i=1}^KP_{e_i}\circ T} < 2^{-k}$.

Fix an index $n \in \{n_1, \ldots, n_K\}$.  Recall that for any $x, v \in \mathcal{H}$, we have the polarization identity
\begin{align*}
4\langle x, v \rangle e_n &= \norm{e_n+v}^2P_{e_n+v}(x) + \norm{e_n-v}^2P_{e_n-v}(x) \\ &+  i\norm{e_n+iv}^2P_{e_n+iv}(x) - i\norm{e_n-iv}^2P_{e_n-iv}(x).
\end{align*}
We also have, for any $x \in \mathcal{H}$,
\[(P_{e_n} \circ T)(x) = \langle T(x), e_n \rangle e_n = \langle x, T^*(e_n)\rangle e_n.\]
Using $v = T^*(e_n)$ in the polarization identity, we thus have
\begin{align*}
P_{e_n} \circ T &= \frac{\norm{e_n+T^*(e_n)}^2}{4}P_{e_n+T^*(e_n)} + \frac{\norm{e_n-T^*(e_n)}^2}{4}P_{e_n-T^*(e_n)}\\&+i\frac{\norm{e_n+iT^*(e_n)}^2}{4}P_{e_n+iT^*(e_n)}-i\frac{\norm{e_n-iT^*(e_n)}^2}{4}P_{e_n-iT^*(e_n)}
\end{align*}
By Lemma \ref{lemma:Adjoint}, $T^*(e_n)$ is a computable vector, so by Corollary \ref{cor:ComputePtProj} each projection operator appearing in the above expression is a computable operator. Thus we have written $P_{e_n}\circ T$ as a linear combination of projections onto computable vectors with computable coefficients.  Summing over $n \in \{n_1, \ldots, n_K\}$ completes the proof.
\end{proof}

We note that the proof of Lemma \ref{lem:FiniteRankApproximation2} is uniform.

\begin{theorem}\label{theorem:EquivalentCompactness}
Suppose that $T$ is a bounded linear operator.  The following are equivalent:
\begin{enumerate}
    \item{$T$ is a computably image-compact operator of $\mathcal{H}^\#$.}
    \item{$T$ is a computable point of $\mathcal{K}(\mathcal{H}^\#)$.}
\end{enumerate}
\end{theorem}
\begin{proof}
The direction (1) implies (2) is immediate from Lemma \ref{lem:FiniteRankApproximation2} and the definition of $\mathcal{K}(\mathcal{H}^\#)$.

For the (2) implies (1) direction, fix $k \in \mathbb{N}$.  Compute a generated point $g$ of $\mathcal{K}(\mathcal{H}^\#)$ so that $\norm{g-T} < 2^{-(k+1)}$.  By definition of the presentation $\mathcal{K}(\mathcal{H}^\#)$, it is possible to compute scalars $\lambda_1, \ldots, \lambda_n$ and $\mathcal{H}^\#$-indices of orthonormal vectors $v_0, \ldots, v_n$ such that $g = \sum_{j=1}^n \lambda_jP_{v_j}$.  Let $X = \operatorname{span}(v_1, \ldots, v_n)$.  It follows that $g[\overline{B}(0_{\mathcal{H}}; 1)] = g[C]$, where $C$ is the closed unit ball of $X$.  By \cite[Proposition 3.7]{McNicholl.2020}, $C$ is computably compact.  Since computable compactness is preserved by computable functions, $\overline{g[\overline{B}(0_{\mathcal{H}}; 1)]}$ is computably compact.

It is thus possible to compute rational open balls $B(x_1; 2^{-(k+1)}), \ldots, B(x_r; 2^{-(k+1)})$ in $X$ so that $\overline{g[\overline{B}(0_{\mathcal{H}}; 1)]} \subseteq \bigcup_{j=1}^rB(x_r; 2^{-(k+1)})$ and so that $B(x_j; 2^{-(k+1)}) \cap g[\overline{B}(0_{\mathcal{H}}; 1)] \neq \emptyset$.  Then $\overline{T[\overline{B}(0_{\mathcal{H}}; 1)]} \subseteq \bigcup_{j=1}^rB(x_j; 2^{-(k+1)})$, so $T$ is computably image-compact.
\end{proof}

Brattka and Dillhage, \cite[Corollary 19]{BrattkaDillhage.2005}, showed that (in our terminology) computable points of $\mathcal{K}(\mathcal{H}^\#)$ have computably compact spectra (this result also follows directly from the combination of our Theorem \ref{theorem:ComputableCompactSpectrumNonunital} and Proposition \ref{prop:PresentationOfCompacts}).  We thus also have the following.

\begin{corollary}
If $T$ is a computably image-compact normal operator of $\mathcal{H}^\#$, then $\sigma(T)$ is computably compact.
\end{corollary}

We conclude with an effective version of the Spectral Theorem for compact normal operators.  To simplify notation, when $\lambda$ is an eigenvalue of an operator we let $P_\lambda$ denote the projection onto the corresponding eigenspace.  In light of Theorem \ref{theorem:EquivalentCompactness} and the existing terminology used in \cite{BrattkaDillhage.2007}, we use the term ``computably compact operator" instead of ``computably image-compact operator".

\begin{theorem}\label{thm:Spectral}
Suppose that $T$ is a computably compact normal operator of $\mathcal{H}^\#$.  Then:
\begin{enumerate}
\item{Every eigenvalue of $T$ is computable.}
\item{From an index of a non-zero eigenvalue $\lambda$ of $T$, it is possible to compute an $\mathcal{H}^\#$-index of $P_\lambda$.}
\item{From $k \in \mathbb{N}$ it is possible to compute an index of a finite set $F$ of non-zero eigenvalues of $T$ so that $\norm{T - \sum_{\lambda \in F}\lambda P_\lambda} < 2^{-k}$.}
\end{enumerate}
\end{theorem}
\begin{proof}
(1): Suppose that $\lambda$ is an eigenvalue of $T$.  If $\lambda = 0$ then $\lambda$ is computable.  Otherwise, $\lambda$ is an isolated point of the computably compact set $\sigma(T)$ (see \cite[Theorem 1.4.11]{Murphy.1990}), and is therefore computable.

(2): Since $\lambda$ is isolated in $\sigma(T) \cup \{0\}$, the characteristic function $\chi_{\lambda}$ of $\{\lambda\}$ is a computable function from $\sigma(T) \cup \{0\}$ to $\mathbb{C}$, so the result follows from Corollary \ref{corollary:ApplyingFuncCalc}, using the presentation $\mathcal{K}(\mathcal{H}^\#)$ and the fact that $P_\lambda = \chi_{\lambda}(T)$.

(3): Being computably compact, $\sigma(a)$ is in particular c.e. closed, and therefore contains a computable dense sequence of points $(\lambda_n)_{n \in \mathbb{N}}$ (see \cite[Lemma 3.27]{DowneyMelnikov}).  Since every point of $\sigma(a)$ (except possibly $0$) is isolated, every non-zero eigenvalue of $T$ occurs (possibly repeatedly) in this dense sequence.  Given a finite set $F \subseteq \mathbb{N}$ we can compute $\norm{T - \sum_{\lambda \in F}\lambda P_\lambda}$ because each $P_\lambda$ is computable, uniformly in $\lambda$ (Corollary \ref{cor:ComputePtProj}), so we can search through finite subsets of $\mathbb{N}$ to find an $F$ such that $\norm{T - \sum_{\lambda \in F}\lambda P_\lambda} < 2^{-k}$.  Since the classical Spectral Theorem tells us that $T = \sum_{n \in \mathbb{N}}\lambda P_\lambda$, this search terminates.
\end{proof}

%\section*{Acknowledgments} \label{sec:ack}

\bibliographystyle{amsplain}
\bibliography{paperbib}


\end{document}
